\documentclass[11pt,a4paper]{article}
\usepackage[margin=22mm]{geometry}
\usepackage{enumitem}
\usepackage{xcolor}
\usepackage{booktabs}
\usepackage{titlesec}
\PassOptionsToPackage{hyphens}{url}
\usepackage[hidelinks]{hyperref}
\sloppy
\emergencystretch=3em
\Urlmuskip=0mu plus 3mu\relax

\definecolor{accent}{HTML}{1F6FB2}
\titleformat{\section}{\large\bfseries\color{accent}}{\thesection}{0.6em}{}
\titleformat{\subsection}{\normalsize\bfseries}{\thesubsection}{0.5em}{}
\setlist[itemize]{leftmargin=1.4em,itemsep=2pt,topsep=3pt}
\setlength{\parindent}{0pt}
\setlength{\parskip}{5pt}

\begin{document}

\begin{center}
  {\LARGE\bfseries NIDS}\\[2pt]
  {\large Network State Detection and Forecasting}\\[6pt]
  {\small Product Description}
\end{center}

\section{Overview}

NIDS is a fully offline network intrusion detection and short-horizon
forecasting system. It moves past per-flow benign/malicious labelling towards
\emph{predictive} defence: it learns how the aggregate state of a network
evolves over time and forecasts the next several state windows, then attaches a
restrained MITRE ATT\&CK interpretation and a per-prediction feature
attribution so a defender can see \emph{why} a call was made.

Everything runs on a local machine. There is no cloud dependency, no telemetry,
and no external API call at inference time.

\section{The problem it addresses}

Traditional classifiers treat each traffic flow in isolation and map it to a
single label. That discards the temporal structure of an intrusion --- the
order in which ports are probed, the way SYN bursts precede floods, the timing
of reconnaissance before lateral movement. An intrusion is a process that
unfolds over time. NIDS models that process: it aggregates traffic into
fixed time windows, learns the state-transition behaviour
$P(S_{t+1}\mid S_t)$, and rolls it forward.

\section{Pipeline}

\begin{enumerate}[leftmargin=1.6em,itemsep=3pt]
  \item \textbf{Capture / ingest.} Live capture via \texttt{dumpcap} (Windows)
  or \texttt{tcpdump} (Linux, macOS), or offline upload of a
  \texttt{.pcap}/\texttt{.pcapng}/\texttt{.csv} file.
  \item \textbf{Flow features.} The Scapy-based \texttt{cicflowmeter} package
  turns packets into the 77-feature CICFlowMeter flow schema.
  \item \textbf{Packet features.} A separate Scapy pass derives packet-level
  attributes the flow schema omits: TTL mean/variance, TCP window size,
  retransmission counts, IP-fragment counts, payload-size distribution.
  \item \textbf{Per-flow state.} A trained Keras ANN (\texttt{ISAA\_ANN.h5},
  MinMax-scaled input) classifies each flow as BENIGN / DDoS / DoS / PortScan.
  \item \textbf{Temporal dataset.} Flows are sorted chronologically, bucketed
  into 10-second windows, and aggregated into a 28-feature network-state
  vector per window, with state transitions and length-5 sliding sequences and
  a strictly chronological train/val/test split.
  \item \textbf{Forecast.} An LSTM predicts the next window's state; a separate
  direct model predicts six windows ahead (H1--H6) with a frozen
  early-warning threshold. Persistence and multinomial-logistic-regression
  baselines are scored on identical sequences.
  \item \textbf{ATT\&CK context.} An offline, evidence-aware rules layer maps
  the predicted state to a small set of ATT\&CK techniques
  (T1595, T1046, T1498, T1499) with a deterministic heuristic confidence,
  supporting evidence, and explicit limitations. It never asserts attribution
  or intent.
  \item \textbf{Explainability.} Gradient\,$\times$\,input attribution over the
  77 scaled features identifies which flow statistics drove each prediction,
  both per flow and aggregated for the capture.
\end{enumerate}

\section{Key features}

\begin{itemize}
  \item Two input modes --- live capture and offline file analysis --- sharing
  one extraction/prediction path.
  \item No fabricated values anywhere: every figure shown is read back from a
  real subprocess or model output, or displayed as \texttt{N/A}.
  \item Chronological-only splits; a dedicated data-leakage check across
  windows and sequences.
  \item One-step and six-step forecasting with honest, labelled proxy-time
  limitations.
  \item Deterministic, offline ATT\&CK mapping with per-candidate evidence and
  limitations.
  \item Per-prediction and per-capture feature attribution.
  \item A Wireshark-inspired local dashboard (static HTML/CSS/JS, no build
  step) served by the backend.
\end{itemize}

\section{Architecture and technology}

\begin{tabular}{@{}ll@{}}
  \toprule
  Backend        & FastAPI + Uvicorn, a small explicit state machine \\
  Model          & TensorFlow~2.18 / Keras Sequential MLP (4-class softmax) \\
  Scaler         & scikit-learn MinMaxScaler (persisted) \\
  Flow features  & \texttt{cicflowmeter} 0.5.0 (Scapy-based) \\
  Packet features& Scapy 2.6.1 streaming \texttt{PcapReader} \\
  Forecasting    & Keras LSTM (one-step and direct H1--H6) \\
  ATT\&CK        & Local pinned Enterprise ATT\&CK v19.1 subset \\
  Frontend       & Plain HTML/CSS/JS, polled REST, no framework \\
  Capture tools  & \texttt{dumpcap}/\texttt{tcpdump}, \texttt{tshark},
                   \texttt{capinfos} \\
  \bottomrule
\end{tabular}

\section{Platform support}

\begin{itemize}
  \item \textbf{Windows} --- install Wireshark; run the PowerShell script
  \texttt{run\_backend.ps1}.
  \item \textbf{Linux} --- \texttt{setcap} on \texttt{tcpdump}; run
  \texttt{run\_backend.sh}.
  \item \textbf{macOS} --- \texttt{scripts/macos\_setup.sh} installs the
  Wireshark CLI tools and loads the ChmodBPF helper (macOS has no
  \texttt{setcap}); the interface picker shows friendly names
  (Wi-Fi, Ethernet, Thunderbolt) mapped from \texttt{networksetup}.
\end{itemize}

Open the dashboard at \texttt{http://127.0.0.1:8765}.

\section{Provenance and attribution}

\subsection*{Base repository}

Cloned from the original \texttt{NIDS} repository (commits by Manigandan
Ramadasan, \texttt{453a5a1}\,--\,\texttt{39316cb}; audited in
\texttt{PIPELINE.md} as of \texttt{main @ 39316cb}). Pre-existing content:
\texttt{models/ISAA\_ANN.h5} (the trained Keras ANN), \texttt{models/minmax.bin}
(its scikit-learn MinMaxScaler), three Jupyter notebooks
(\texttt{NIDS -- ANN.ipynb}, \texttt{NIDS -- ML.ipynb},
\texttt{NIDS (Prediction).ipynb}), \texttt{scripts/packetsniff.sh}, and a report
PDF. The prediction notebook's hardcoded \texttt{O:\textbackslash VIT\ldots} and
\texttt{C:\textbackslash Users\textbackslash varun\ldots} paths indicate a
Vellore Institute of Technology student project; ``ISAA\_ANN'' is the original
team's model name. Everything from commit \texttt{945f9ad} onward --- the
end-to-end pipeline, FastAPI backend, dashboard, Phase 3/4 LSTM forecasting,
MITRE layer, gradient\,$\times$\,input explainability, packet-level features, and
cross-platform capture --- was added on top.

\subsection*{Feature extraction}

\texttt{cicflowmeter} 0.5.0 (PyPI, MIT-licensed, Scapy-based) --- a drop-in for
the Java \texttt{ahlashkari/CICFlowMeter}, which was abandoned here because its
\texttt{jnetpcap} dependency has no working modern 64-bit build path. Two
upstream bugs are patched idempotently by \texttt{scripts/patch\_cicflowmeter.py}.

\subsection*{Training and evaluation dataset}

CICIDS2017, produced by the Canadian Institute for Cybersecurity, University of
New Brunswick --- \url{https://www.unb.ca/cic/datasets/ids-2017.html}. Eight
official weekday session flow CSVs are used: Phase 3 trains on Monday, Tuesday,
Wednesday and Thursday-Afternoon-Infiltration; Phase 4 adds Thursday-morning Web
Attacks, Friday morning, Friday PortScan and Friday DDoS. The raw CSVs are not
vendored in Git --- they are published as GitHub Release assets
(\texttt{phase3-datasets},
\texttt{nids-cicids2017-phase4-added-source-csvs.tar.gz}) with per-file SHA-256
hashes in \texttt{datasets/SHA256SUMS} and the archive hash in
\texttt{datasets/README.md}. The ANN was trained by the original team on the
classic CICIDS2017 78-column export (77 numeric features after dropping the
duplicate header column and the label).

\subsection*{MITRE ATT\&CK}

Enterprise ATT\&CK v19.1 (MITRE Corporation), distributed as STIX~2.1 via
\url{https://github.com/mitre-attack/attack-stix-data}. A pinned, verified subset
--- techniques T1595, T1046, T1498, T1499; tactics Reconnaissance, Discovery,
Impact --- is committed at
\texttt{backend/mitre/data/enterprise\_attack\_v19\_1\_subset.json} and loaded
with publisher and bundle-SHA-256 provenance checks
(\texttt{backend/mitre/mapper.py}). Used entirely offline; no network calls.

\subsection*{Demonstration captures}

\texttt{captures.zip} is the community ``DDoS Packet Capture Collection'' (mirror
\url{https://wqrld.net/captures.zip}) --- anonymized real-world and lab DDoS
captures, free to use for building DDoS-protection systems, with attribution
appreciated but not required. Members are named \texttt{amp.*} (amplification) or
\texttt{pkt.*} (spoofed / direct). It is not part of the repository; four small
members are copied into \texttt{pcaps/} locally for demos.

\subsection*{Core libraries}

FastAPI and Uvicorn, TensorFlow~2.18 / Keras, scikit-learn~1.5.2, Scapy~2.6.1,
pandas / NumPy, joblib, python-multipart, and pytest --- all pinned in
\texttt{backend/requirements.txt}.

\section{Known limitations}

\begin{itemize}
  \item The ANN was trained on CICIDS2017 (2017 lab traffic) and can
  misclassify modern HTTPS/QUIC/STUN as DoS/DDoS --- a model-generalisation
  limit, not a pipeline defect.
  \item Only the ANN is deployable; Random Forest / Decision Tree / KNN from
  the original notebooks were never persisted.
  \item Forecast horizons use row-order proxy windows on the training CSVs,
  not validated wall-clock seconds; the UI labels this explicitly.
  \item The ATT\&CK layer covers Reconnaissance / Discovery / Impact only;
  it does not cover Initial Access, Lateral Movement, C2 or Exfiltration
  (the classifier has no such states).
  \item The packet-feature extractor is standalone --- this
  \texttt{cicflowmeter} build emits no source/destination IP columns, so the
  two feature levels are not yet joined on a common flow key.
  \item The persisted scaler was pickled under an older scikit-learn; values
  are unaffected and the version warning is suppressed at load.
\end{itemize}

\end{document}
